\documentclass[11pt,a4paper]{article}
\usepackage[margin=1in]{geometry}
\usepackage{amsmath,amssymb,amsthm}
\usepackage{algorithm}
\usepackage{algpseudocode}
\usepackage{booktabs}
\usepackage{hyperref}

\theoremstyle{definition}
\newtheorem{definition}{\textbf{Definition}}[section]
\newtheorem{theorem}{\textbf{Theorem}}[section]
\newtheorem{lemma}[theorem]{\textbf{Lemma}}
\newtheorem{remark}{\textbf{Remark}}[section]

\title{\textbf{Gradient Accumulation: Simulating Virtual Batch Sizes Under Memory Constraints}}
\author{\textbf{Team Scaibu} \\ \small Email: \texttt{jaydeep.9664920749@gmail.com} \\ \small \textbf{Architecture and Core Engine Team}}
\date{\textbf{October 2026}}

\begin{document}
\maketitle

\section{Definitions and Cognitive Mental Models}

\subsection{Formal Mathematical Definition}
\textbf{Definition (Gradient Accumulation and Normalized Microbatching):}
Optimal generalization and computational stability in large language model pre-training require massive global batch sizes $B_{\text{target}} \in [10^6, 8 \times 10^6]$ tokens (Kaplan et al., 2020; McCandlish et al., 2018). However, the physical GPU memory capacity bounds the maximum batch size that can be processed concurrently in a single forward-backward pass to a smaller microbatch size $b \ll B_{\text{target}}$. \textbf{Gradient Accumulation} overcomes this by summing gradients across $K$ consecutive microbatches before invoking an optimizer step:
{\boldmath
\begin{equation}
g_{\text{accum}} = \sum_{k=1}^K g_k, \quad g_k = \sum_{x \in \mathcal{B}_k} \nabla_\theta \ell(x; \theta)
\end{equation}
}
where the accumulation factor satisfies $K = B_{\text{target}} / (b \cdot N_{\text{devices}})$.

To ensure mathematical equivalence with a true monolithic batch of size $B_{\text{target}}$, the accumulated gradient must be appropriately normalized:
\begin{enumerate}
    \item \textbf{Microbatch Step Averaging (Uniform Lengths)}:
    {\boldmath
    \begin{equation}
    g_{\text{opt}} = \frac{1}{K} g_{\text{accum}} = \frac{1}{K} \sum_{k=1}^K g_k
    \end{equation}
    }
    \item \textbf{Token-Weighted Averaging (Variable Document Lengths)}:
    When sequence packing (#230) or unpadded batches have varying token counts $T_k = |\mathcal{B}_k|$, dividing by $K$ introduces token-weight distortion. Proper token-level normalization weights each microbatch by its token fraction:
    {\boldmath
    \begin{equation}
    g_{\text{opt}} = \frac{1}{\sum_{k=1}^K T_k} \sum_{k=1}^K g_k \cdot T_k
    \end{equation}
    }
\end{enumerate}
In distributed data-parallel settings, network collective synchronizations (All-Reduce, #233) are bypassed during microbatches $1 \dots K-1$ and triggered exclusively on microbatch $K$, reducing network collective overhead by factor $K$.

\subsection{Expanded Non-Technical Definition (Intuition for Anyone)}
\begin{enumerate}
    \item \textbf{Paying with Coins in a Piggy Bank}: If you want to pay a \$100 bill ($B_{\text{target}}$) but your wallet only holds \$10 at a time ($b$), you don't run to the store 10 times to buy 10 separate tiny items. You drop \$10 into a piggy bank 10 times ($K=10$), and then pay the entire \$100 all at once!
    \item \textbf{Simulating a Supercomputer on a Laptop}: Even with a single consumer graphics card that can only fit 2 sentences in memory, you can accumulate gradients 1,000 times to simulate the training dynamics of a 2,000-sentence supercomputer run.
    \item \textbf{Saving Network Bandwidth}: In a multi-GPU cluster, GPUs only talk to each other when the piggy bank is full, keeping the network completely quiet during the accumulation steps.
\end{enumerate}

\subsection{Child-Friendly Mental Model (Physical Metaphor)}
Imagine filling a large water barrel with a coffee mug:
\begin{itemize}
    \item \textbf{Naive Update}: Pouring one mug into the barrel, walking 100 feet to turn on the main valve, walking back, pouring another mug, walking 100 feet... You spend all your time walking!
    \item \textbf{Gradient Accumulation}: Standing still and dumping 20 mugs into the barrel ($K=20$), and only walking to the valve once the barrel is completely full.
\end{itemize}

\section{Core Mathematical Equations and Definitions}

\subsection{Effective Global Batch Size}
{\boldmath
\begin{equation}
B_{\text{global}} = K \cdot N_{\text{data-parallel}} \cdot b_{\text{micro}}
\end{equation}
}
\begin{itemize}
    \item \textbf{Formal Definition}: Total effective batch size synthesized per optimizer weight update.
    \item \textbf{What It Means}: Enables scaling batch size independently of physical GPU count.
    \item \textbf{Why It Is Important}: Allows small clusters to match the critical batch size (#244) of massive clusters.
\end{itemize}

\subsection{Communication Reduction Factor}
{\boldmath
\begin{equation}
\mathcal{R}_{\text{comm}} = \frac{1}{K}
\end{equation}
}
\begin{itemize}
    \item \textbf{Formal Definition}: Network All-Reduce frequency under gradient accumulation.
    \item \textbf{What It Means}: Suppresses inter-device communication by factor $K$.
    \item \textbf{Why It Is Important}: Massive acceleration on clusters with slow inter-node interconnects.
\end{itemize}

\section{Rigorous Mathematical Analysis Checklist}

\begin{itemize}
    \item \textbf{Notation}: $K$ (accumulation steps), $b$ (microbatch size), $g_k$ (microbatch gradient), $g_{\text{accum}}$ (accumulated sum).
    \item \textbf{Epistemic Status}: Standardized optimization technique supported in PyTorch, JAX, and DeepSpeed.
    \item \textbf{Dependency Graph}: Microbatch $k \to \text{Backward} \to \text{Accumulate } g_k \to \text{Check } k == K \to \text{Normalize and Step}$.
    \item \textbf{Dimensions}: Model parameter dimension $P$.
    \item \textbf{Regimes}: Batch size scaling under memory constraints, large language model pre-training.
    \item \textbf{Invariants}: In the absence of stateful batch normalization, output parameter updates are identical to monolithic batch updates.
    \item \textbf{Isomorphisms}: Isomorphic to delayed numerical integration in ODE solvers.
\end{itemize}

\section{Theoretical Theorems and Formal Proofs}

\begin{theorem}[Linear Linearity of Stochastic Gradient Accumulation]
Let loss $\mathcal{L}(\theta) = \frac{1}{B} \sum_{i=1}^B \ell_i(\theta)$ be the empirical risk over a monolithic batch of size $B = K \cdot b$. The normalized accumulated gradient $\bar{g}_{\text{accum}} = \frac{1}{K} \sum_{k=1}^K \left( \frac{1}{b} \sum_{j \in \mathcal{B}_k} \nabla \ell_j(\theta) \right)$ is identically equal to the monolithic batch gradient $\nabla \mathcal{L}(\theta)$:
{\boldmath
\begin{equation}
\bar{g}_{\text{accum}} \equiv \nabla \mathcal{L}(\theta)
\end{equation}
}
\end{theorem}

\begin{proof}
Let the global batch $\mathcal{B} = \{1, \dots, B\}$ be partitioned into $K$ disjoint microbatches $\mathcal{B}_1, \dots, \mathcal{B}_K$, each of size $b$, such that $\bigcup_{k=1}^K \mathcal{B}_k = \mathcal{B}$ and $\mathcal{B}_k \cap \mathcal{B}_l = \emptyset$ for $k \neq l$.
The gradient of the monolithic batch loss is:
{\boldmath
\begin{equation}
\nabla \mathcal{L}(\theta) = \nabla_\theta \left( \frac{1}{B} \sum_{i=1}^B \ell_i(\theta) \right) = \frac{1}{B} \sum_{i=1}^B \nabla_\theta \ell_i(\theta)
\end{equation}
}
By partitioning the summation over the $K$ disjoint subsets:
{\boldmath
\begin{equation}
\frac{1}{B} \sum_{i=1}^B \nabla_\theta \ell_i(\theta) = \frac{1}{B} \sum_{k=1}^K \sum_{j \in \mathcal{B}_k} \nabla_\theta \ell_j(\theta)
\end{equation}
}
Substituting $B = K \cdot b$:
{\boldmath
\begin{equation}
\frac{1}{K \cdot b} \sum_{k=1}^K \sum_{j \in \mathcal{B}_k} \nabla_\theta \ell_j(\theta) = \frac{1}{K} \sum_{k=1}^K \left( \frac{1}{b} \sum_{j \in \mathcal{B}_k} \nabla_\theta \ell_j(\theta) \right) = \frac{1}{K} \sum_{k=1}^K g_k = \bar{g}_{\text{accum}}
\end{equation}
}
Because differentiation and finite summation commute linearly, the accumulated gradient is mathematically identical to the monolithic batch gradient.
\end{proof}

\begin{theorem}[Token-Weighted Invariance Under Variable-Length Microbatches]
If microbatches possess unequal token counts $T_k$, the token-weighted estimator $\hat{g} = \frac{\sum_{k=1}^K g_k}{\sum_{k=1}^K T_k}$ is an unbiased estimator of the global token-level loss gradient, whereas uniform step averaging $\frac{1}{K} \sum_{k=1}^K \frac{g_k}{T_k}$ introduces an asymptotic bias bounded by $\mathcal{O}(\operatorname{Cov}(T_k, g_k / T_k))$.
\end{theorem}

\begin{proof}
Let $T_{\text{total}} = \sum_{k=1}^K T_k$ be the total tokens. The token-level loss is:
{\boldmath
\begin{equation}
\mathcal{L}_{\text{token}} = \frac{1}{T_{\text{total}}} \sum_{k=1}^K \sum_{t=1}^{T_k} \ell_{k, t} \implies \nabla \mathcal{L}_{\text{token}} = \frac{1}{T_{\text{total}}} \sum_{k=1}^K g_k
\end{equation}
}
where $g_k = \sum_{t=1}^{T_k} \nabla \ell_{k, t}$.
Under uniform step averaging, each microbatch is first divided by its local length $T_k$: $\bar{g}_k = g_k / T_k$, and then averaged over $K$:
{\boldmath
\begin{equation}
g_{\text{step}} = \frac{1}{K} \sum_{k=1}^K \bar{g}_k
\end{equation}
}
Notice that:
{\boldmath
\begin{equation}
\mathbb{E}[g_{\text{step}}] = \mathbb{E}\left[ \frac{1}{K} \sum_{k=1}^K \bar{g}_k \right] = \mathbb{E}[\bar{g}_k]
\end{equation}
}
Whereas the true token gradient expectation is:
{\boldmath
\begin{equation}
\mathbb{E}[\nabla \mathcal{L}_{\text{token}}] = \frac{\mathbb{E}[T_k \bar{g}_k]}{\mathbb{E}[T_k]} = \mathbb{E}[\bar{g}_k] + \frac{\operatorname{Cov}(T_k, \bar{g}_k)}{\mathbb{E}[T_k]}
\end{equation}
}
If longer sequences have systematically higher or lower gradient magnitudes (e.g. attention perplexity), $\operatorname{Cov}(T_k, \bar{g}_k) \neq 0$, causing uniform step averaging to produce biased gradients. Token-weighted normalization eliminates this covariance term identically.
\end{proof}

\section{Foundational Architectural Questions and Epistemic Meta-Questions}

\subsection{Architectural Question 1: Interaction with Batch Normalization}
\textbf{Question:} Why does gradient accumulation behave differently in architectures with BatchNorm vs. LayerNorm/RMSNorm?\\
\textbf{Answer:} 
\begin{itemize}
    \item \textbf{LayerNorm / RMSNorm (Transformers)}: Normalizes across the hidden dimension of each token independently. Microbatch size $b$ has zero effect on normalization statistics, ensuring exact bit-level equivalence with large batches.
    \item \textbf{BatchNorm (CNNs)}: Calculates mean and variance across the batch dimension. Accumulating across microbatches computes normalization over small $b$, producing noisier statistics than a monolithic batch $B$.
\end{itemize}

\subsection{Architectural Question 2: Learning Rate and Warmup Scaling}
\textbf{Question:} Should the learning rate be scaled when increasing accumulation steps $K$?\\
\textbf{Answer:} Yes. Increasing $K$ from 1 to 8 scales the effective global batch size by $8\times$. Following the linear scaling rule (#233, Goyal et al.), the base learning rate should be scaled proportionally ($8\times$) or using the square root scaling rule ($\sqrt{8}\times$), accompanied by an extended learning rate warmup period to prevent early optimization divergence.

\section{Algorithmic Implementation Specification}

\begin{algorithm}[H]
\caption{Gradient Accumulation and Step Gating}
\begin{algorithmic}[1]
\Require Accumulated buffer $\mathbf{g}_{\text{acc}} \in \mathbb{R}^P$, microbatch grad $\mathbf{g}_{\text{mb}} \in \mathbb{R}^P$, step $k$, total $K$, mode, tokens $T_{\text{mb}}, T_{\text{tot}}$
\Ensure Updated buffer $\mathbf{g}_{\text{acc}}'$, is\_update flag, step grad $\mathbf{g}_{\text{opt}}$, count $k$
\State Allocate $\mathbf{g}_{\text{acc}}' \in \mathbb{R}^P$
\If{$\text{mode} == \text{"token"}$}
    \State $w \gets T_{\text{mb}} / \max(1, T_{\text{tot}})$
    \For{$i = 0$ \textbf{to} $P-1$}
        \State $\mathbf{g}_{\text{acc}}'[i] \gets \mathbf{g}_{\text{acc}}[i] + \mathbf{g}_{\text{mb}}[i] \cdot w$
    \EndFor
\Else
    \For{$i = 0$ \textbf{to} $P-1$}
        \State $\mathbf{g}_{\text{acc}}'[i] \gets \mathbf{g}_{\text{acc}}[i] + \mathbf{g}_{\text{mb}}[i]$
    \EndFor
\EndIf
\State $\text{is\_update} \gets (k == K)$
\If{$\text{is\_update}$}
    \If{$\text{mode} == \text{"step"}$}
        \State $\mathbf{g}_{\text{opt}} \gets (1.0 / K) \cdot \mathbf{g}_{\text{acc}}'$
    \Else
        \State $\mathbf{g}_{\text{opt}} \gets \mathbf{g}_{\text{acc}}'$
    \EndIf
\Else
    \State $\mathbf{g}_{\text{opt}} \gets [\,]$
\EndIf
\State \Return $(\mathbf{g}_{\text{acc}}', \text{is\_update}, \mathbf{g}_{\text{opt}}, k)$
\end{algorithmic}
\end{algorithm}

\section{Big-$\Theta$ Complexity Analysis}

\begin{itemize}
    \item \textbf{Vector Accumulation Time}: $\Theta(P)$ additions per microbatch.
    \item \textbf{Step Normalization Time}: $\Theta(P)$ scaling operations on step $K$.
    \item \textbf{Total Computational Overhead}: $\Theta(P)$ linear operations.
    \item \textbf{Memory Storage Overhead}: Exactly $\Theta(P)$ floats for the accumulation buffer.
\end{itemize}

\section{Single Canonical Repository Reference}

The complete deterministic scalar implementation, verification suite, and unit test benchmarks are published at:
\begin{center}
\url{https://github.com/Chief-Strategist-J/llm-obs-infra}
\end{center}

\end{document}
